\subsection{無限遠の一点を持つ加群と二次行列因子分解}
\label{subsec:cubic-CY-three-alpha-elliptic-line-modules}

\begin{theorem}\label{thm:cubic-CY-three-alpha-elliptic-filtered-line-modules}
滑らかな楕円境界の平行移動は無限位数とする。
アフィン包絡 $D$ 上のfiltered巡回左加群のうち、
随伴加群が一点 $P\in E$ に対応するpoint moduleとなるものを考える。
$V=S_1$ とし、$L\subset V$ を $P$ を消す二次元空間とする。
$L$ の基底を $\ell_1,\ell_2$ とし、先頭の唯一の一次syzygyを
\[
 m_1\ell_1+m_2\ell_2=\sum_i k_i r_i
\]
と書く。ここで $r_i$ は三つの斉次二次関係である。
包絡におけるその右辺を $T_1+T_0$ と書く。
$T_1\in V$、$T_0\in k$ は低次係数から線型に求まる。

全filtered持ち上げは、二つのスカラー $\alpha_i$ で作る
$\ell'_i=\ell_i+\alpha_i$ により、次の式で記述できる。
\[
 T_1+\alpha_1m_1+\alpha_2m_2+c_1\ell_1+c_2\ell_2=0,
 \qquad T_0+c_1\alpha_1+c_2\alpha_2=0 .
\]
第一式は $(\alpha_1,\alpha_2)$ に対するアフィン直線を与え、$c_1,c_2$ を一意に定める。
第二式はその直線に二次以下の条件を課す。
得られる加群
\[
 M=D/(D\ell'_1+D\ell'_2)
\]
は $\operatorname{End}_D(M)=k$ を満たす。
従って中央三次 $\omega$ は一つの値 $\lambda$ で作用し、
$M$ は $U_\lambda$ の無限次元加群となる。
\end{theorem}
\begin{proof}
左point moduleの最小解消は
\[
 0\longrightarrow S(-2)\longrightarrow S(-1)^2
 \longrightarrow S\longrightarrow M_P\longrightarrow0
\]
である。実際、点関係の行列は滑らかな点曲線上で階数2であり、
$V\otimes L$ 内に属する二次関係は一次元である。
その関係の左二係数 $m_1,m_2$ は、整域性から独立である。
二本で割った巡回加群は曲線から得る点加群へ全射する。
前の次数が一次元なら、点関係をその生成元に掛けると、
右端の一次元評価により次の次数に二つの独立な一次制約を得る。
点グラフ上の全評価で階数は2であるため、帰納的に各次数の次元は1である。
従って二本で点加群の全関係を生成し、Hilbert級数の比較で
上記の一次syzygyが全核を生成する。

随伴を $M_P$ とするfiltered加群では、元 $\ell'_1,\ell'_2$ は
最高重みを下げる帰納法によりその巡回関係を生成する。
同じ帰納法で、先頭syzygyは $(m_1+c_1,m_2+c_2)$ の形に持ち上がる。
二つの積を展開し、次数1と0を分けると表示した二式になる。
逆に二式が成立すれば、持ち上げた複体の随伴は上の点加群解消なので、
最高重みの降下により完全性と所望の随伴が従う。
$m_1,m_2$ の張る二平面の基点を求める。
左点加群の点列の進み方を $\tau$ と書く。
関係を $(Q,\tau^{-1}Q)$ で評価すると、
$E$ 上で $[m_1:m_2]$ の有理写像は
$[-\ell_2:\ell_1]\circ\tau^{-1}$ と一致する。
$\ell_i$ の鉛筆の基点因子は $P$ で、基点を除いた写像の次数は2である。
従って $m_i$ の鉛筆も基点を一つ持ち、その点 $P'$ は
$\mathcal L(-P')\simeq(\tau^{-1})^*\mathcal L(-P)$ で与えられる。
平行移動の量を $s$ とすると、次数2の線型束に対する逆平行移動の引き戻しは
Picard座標に $2s$ を加えるので、$P'=\tau^{-2}(P)$ である。
従って無限位数ではこの平面は $L$ と異なる。
従って $V/L$ に射影した第一式の係数は非零であり、
残りの二成分は $c_1,c_2$ を一意に定める。

$D$ 加群準同型 $M\to M$ は巡回生成元の像で決まる。
その像が最低フィルター次数 $j$ にあれば、重み $j$ の随伴準同型を誘導する。
$j>0$ の非零な随伴準同型は、点加群の次数0生成元を
$j$ 次尾部に送り、その消滅一次平面を一致させる。
しかし尾部の点は $P$ の平行移動の $j$ 回反復であり、$P$ に等しくない。
従って $j=0$ で、準同型はスカラーに限られる。
中央元の作用も準同型なので、特定の切断へ降りる。
\end{proof}

\begin{proposition}\label{prop:cubic-CY-three-alpha-elliptic-line-factorization}
前定理の $M$ と切断 $U_\lambda$ について、
\[
 \operatorname{pd}_{U_\lambda} M\in\{1,\infty\}.
\]
無限射影次元は $\operatorname{Ext}^4_{U_\lambda}(M,M)\ne0$ と同値である。
さらにこの判定に使う二周期解消は、$2\times2$ の行列因子分解として書ける。
\end{proposition}
\begin{proof}
$f=\omega-\lambda$ とし、左自由加群を行ベクトルで書く。
解消の微分を $\ell=(\ell'_1,\ell'_2)^{\rm t}$ と
$m=(m_1+c_1,m_2+c_2)$ とする。
$f$ が巡回生成元を消すので、$a\ell=f$ となる行 $a$ が存在する。
フィルター付き完全性により、その成分の次数は高々2に取れる。
さらに
$(fI_2-\ell a)\ell=0$ であるため、
$b m=fI_2-\ell a$ となる列 $b$ が存在し、その成分も次数高々2となる。
$m\ell=0$ と微分の単射性から $mb=f$、$ab=0$ を得る。
従って
\[
 P=\begin{pmatrix}m\\a\end{pmatrix},\qquad Q=(b\ \ell)
 \qquad\text{は}\qquad PQ=QP=fI_2
\]
を満たす。$f$ の正則性から、商上で交互に $P,Q$ を掛ける二周期複体は完全である。
例えば $vP$ が $f$ の倍数なら $vP=u f$ と持ち上げて $Q$ を掛け、
$f$ で割ると $v=uQ$ となる。反対の偶奇も同様である。
$v\ell=0$ が商で成立すれば $D$ 上で $v\ell=uf=ua\ell$ と持ち上がり、
$v-ua$ は $m$ の倍数となる。従って最後を $\ell$ に接続して
$M$ の射影解消を得る。

$D$ 上の長さ2の解消と中央正則商の変換により、
$\operatorname{Ext}^{j}_{U_\lambda}(M,U_\lambda)=0$ が $j>1$ で成り立つ。
有限射影次元なら、最後の双対微分の分裂により射影次元は高々1となる。
$M$ は非零イデアルによる整域の商で、自由加群の直和因子にはなれない。
従って有限の場合の射影次元は1である。

$U_\lambda\otimes_D^{\mathbf L} M$ は完全で、
中央正則式の変換射 $M\to M[2]$ は完全複体を割った導来商で同型となる。
$M$ が完全でなければその二乗は導来商でも非零なので、
四次自己Extも非零となる。有限の場合の四次Extは零である。
\end{proof}

\begin{corollary}\label{cor:cubic-CY-three-alpha-elliptic-line-finite-pd-test}
前命題の四次自己Extは、表示した二周期解消のコチェインを
内部重み高々 $-3$ に切った複体で正確に計算できる。
次数3,4,5の三項の次元は $4,7,10$ である。
従って、この形の無限次元加群がH3を排除する条件は、
有限二行列の階数の和が7より小さいことに等しい。
\end{corollary}
\begin{proof}
二つの行列因子の最高重みを取ると、$g$ に対する同じ構成の点加群解消になる。
その双対コチェインの随伴複体は $C=S/(g)$ 上の点加群のものとなる。
まず $S$ 上の長さ2の解消を使う。
内部重み $j$ における点加群のコチェインは
$(M_P)_j,(M_P)_{j+1}^{\oplus2},(M_P)_{j+2}$ である。
第一微分は $P_j=\tau^jP$ での $\ell_i$ の評価で、$j=0$ のときだけ
非零のsourceを持ちつつ零になる。
第二微分は $P_{j+1}$ での $m_i$ の評価で、
非零のsourceがある $j\ge-1$ では常に非零である。
実際、$m_i$ の基点は $P_{-2}$ である。
従って絶対自己Extは、次数0の重み0、次数1の重み $-1,0$、
次数2の重み $-2$ にそれぞれ一次元を持ち、他の成分は零である。

中央三次商への変換完全列では、演算子 $h$ は内部重み $-3$ で、
絶対から相対へ戻る写像は重み3である。
次数1と2から戻る写像は上記重みからいずれも零となる。
実際、次数1からは相対次数0の重み正へ、
次数2からは相対次数1の重み1へ移る。
従って相対次数3は $h$ を掛けた次数1で重み $-4,-3$、
相対次数4は $h$ を掛けた次数2で重み $-6,-5$ にある。
特に次数3,4ともに重み $-3$ より上にコホモロジーがない。
最高重みの降下により、次数4の閉元は重み $-3$ 以下に取れ、
その範囲で境界となる閉元の原始元も同じ範囲に取れる。
これで切った複体が正確に第四自己Extを計算する。
対応する自由項の重みは、次数3が4を二本、次数4が5と6、
次数5が7を二本である。
加群の次数 $d$ 以下は次元 $d+1$ なので、主張の三つの次元と階数条件を得る。
\end{proof}

\begin{proposition}\label{prop:cubic-CY-three-alpha-elliptic-line-critical-value}
包絡 $D$ を固定し、上で得たfiltered加群の係数schemeを $\mathcal M_D$ とする。
中心値写像を $\lambda_D:\mathcal M_D\to\mathbb A^1$ と書く。
点 $M$ で $d\lambda_D$ が零であることと、$M$ の切断上での射影次元が
無限であることは同値である。
\end{proposition}
\begin{proof}
$D$ 上の長さ2の解消に $\operatorname{Hom}(-,M)$ を取る。
先頭の計算で一次Extは重み $-1,0$ にしかなく、零次Extに正重みはない。
従って一次Extは重み0以下の閉元で計算でき、境界も同じ範囲で取れる。
その範囲の境界は零である。実際、零次コチェインのスカラー生成元は
$\ell'_1,\ell'_2$ に消される。
よって一次Extは、$M_{\le1}$ の二元で、$m$ により消されるものからなる。

これは $\mathcal M_D$ の接空間に等しい。
二元を $D_{\le1}$ に持ち上げて $\ell'_i$ の一次変分とする。
閉元条件により、元のsyzygyとの積の変分は巡回イデアルに入り、
strictな表示を使って次数1以下の係数で $\ell'_i$ の組合せに書ける。
従ってsyzygyも同時に一次変形でき、双対数上のfiltered点加群を得る。
その先頭一次平面は二次関係の点行列の行列式を消すので、相対点曲線上にある。
持ち上げの差は二つの関係の基底変更である。
逆に関係の一次変形から同じ閉元が得られる。

中央正則商の変換列には
\[
 \operatorname{Ext}^1_D(M,M)\longrightarrow
 \operatorname{End}_{U_\lambda}(M)\xrightarrow{\,h\,}
 \operatorname{Ext}^2_{U_\lambda}(M,M)
\]
がある。最初の射は、$D$ 上の自己拡大における
$\omega-\lambda$ の上三角成分 $M\to M$ を取る写像である。
実際、中央商の二項解消を使うと、この成分が変換射の境界になる。
従ってこれは中心値の微分である。
もし微分が非零なら、右の $\operatorname{End}=k$ に全射して $h(1)=0$ となり、
特異性圏での周期同型から $M$ は完全である。
逆に完全なら射影次元は1なので二次Extは零、微分は全射である。
\end{proof}

\begin{theorem}\label{thm:cubic-CY-three-alpha-elliptic-diagonal-classification-123}
先頭 $(p,q,r)=(1,2,3)$ を固定する。
全低次係数および全切断値について、八次判別式が消えるなら、
その切断には上の形の無限射影次元の加群が存在する。
従ってこの固定境界でH3は八次判別式の補集合に厳密に一致する。
\end{theorem}
\begin{proof}
まず係数全体で $\mathcal M_D$ を動かし、包絡と切断を合わせた係数空間への
有限型射 $\mathcal M\to\{(D,\lambda)\}$ を考える。この射はproperである。
実際、DVR上で一般点の加群をRees同次化し、その次数付き巡回イデアルを
整数係数の切断環と交叉させる。
各次数の商はtorsion-freeで、階数 $n+1$ の自由加群となる。
特殊化した巡回加群を $N$ とする。
もし $t$ がその次数1関係に入れば、境界座標環 $C$ 上の
非零一次関係を持つ巡回商になる。
$C$ は整域で、その一次関係一つによる商は次数ごとに高々3次元なので矛盾する。
従って次数1の二関係は $C_1$ に独立に射影される。
$N/tN$ はその二一次関係による境界の巡回商なので各次数高々一次元である。
これは点行列の階数が全点で2以上であることと、点加群の次数帰納から従う。
一方、$N_n$ は次元 $n+1$ だから、この差は各次数で少なくとも1である。
従って $t$ は $N$ に正則で、$N/tN$ はpoint moduleになる。
次数1の関係部分束はこの特殊化で開チャート内にあり、
そこで上述の係数方程式がDVR上でも成立する。
従って加群の係数点が延長する。一意性は次数1部分束の分離性から従い、
付値判定でproper性を得る。

無限射影次元の加群は、系\ref{cor:cubic-CY-three-alpha-elliptic-line-finite-pd-test}
の有限行列の階数条件で閉部分をなす。
行列因子の係数は有限次数の自由束間のstrictな表示で局所的に選べるため、
同じ階数条件が係数scheme上で使える。
従ってその切断係数空間への像は閉である。
この像は八次判別式の零点に含まれる。

次に一般の固定包絡 $D$ で、少なくとも一つの有限な臨界値が存在することを示す。
点 $P$ を動かすと、第一式の解は $E$ 上のアフィン直線束であり、
第二式はこの直線に沿う二次式である。
$L$ の行列式の次数は $-3$。右評価の像は次の点の二平面に
$\mathcal O_E(1)$ を掛けた束なので次数3であり、
従ってsyzygy直線 $J$ の次数は $-3$ である。
アフィン直線束の方向束 $K$ は
$L^\vee\to J^\vee\otimes\mathcal O_E(1)$ の核なので、次数は $-3$ である。
二次式の最高係数は $J^\vee\otimes K^{-2}$ の切断であり、次数は9である。
これは零切断ではない。
実際、変曲点 $(1:-1:0)$ では
$\ell_1=x+y,\ell_2=z$ と取れ、
$m_1=3x+6y+2z,\ m_2=4x+y+9z$、
同次方向 $\alpha_1=\alpha_2=u$ で最高項は $-18u^2$ になる。

従って最高係数には零点がある。
アフィン直線束を相対射影直線束に完備化すると、
その零点で二次方程式の無限遠の根を得る。
一般の定数係数では、そのfiberで二次式が恒等的に零とはならない。
これはsyzygyの定数項に定数係数の非零線型汎関数 $T_0$ が現れるためである。
よって無限遠の根は $E$ を支配する成分に属する。
その枝では点が $\mathcal M_D$ から脱出するので、
proper性により中心値は有限なままではいられない。
従って該当する成分上で $\lambda_D$ は非定数である。
この成分の射影正規化は $E$ に支配的に写るので種数は1以上である。
全包絡を動かしたaffineな線加群schemeは滑らかである。
第一式は非零の線型係数で解け、第二式は定数係数の一つについて
非零の一次式なので、曲線上のアフィン束から局所的に作れる。
従って標数零で一般の包絡上の $\mathcal M_D$ は滑らかとなる。
正規化した曲線から $\mathbb P^1$ への非定数写像は、
Riemann--Hurwitzにより無限遠の外にも分岐を持つ。
実際、次数 $d$ と種数 $g\ge1$ に対し全分岐量は $2g-2+2d$、
無限遠上の寄与は高々 $d-1$ である。
proper性により有限分岐点は $\mathcal M_D$ に属する。
そこで微分は零になり、前命題が無限射影次元の加群を与える。

幾何学的一般の係数では、この無限射影次元の加群の中心値は
八次判別式の一根である。
定理\ref{thm:cubic-CY-three-alpha-elliptic-eight-monodromy}により
その八根は係数体のGalois群で推移的に移るので、全根がこの像に入る。
従って閉な像は幾何学的に既約な判別式の全体である。
判別式の外は既に双加群滑らかでH3なので、主張を得る。
\end{proof}
